\chapter{Local and Collective Free Energy}
\label{ch:local-collective-elbo}

The bundle geometry of \Cref{ch:geometry} is the inherited foundation of this
chapter.  The probabilistic construction below is additional structure.  It is
not inferred from the existence of the principal bundle, its group, or its two
associated statistical fibers.  The construction answers the local--global
question by putting every interaction record in one fixed normalized joint law
and then deriving each local objective by disintegration of that law.  In this
sense local and collective free energies are compatible, but they are not the
same accounting object: the graph-local objectives below are coordinate
potentials, and summing their incident-factor formulas overcounts shared
records.

\section{A normalized interaction-record model}
\label{sec:local-interaction-record-model}

\hypothesisheading{Agent state and interaction kernels}{hyp:local-interaction-kernels}
Let $(\mathsf X,\mathscr X)$ be a standard Borel context space, let $V$ be an
arbitrary finite set of agents, and let $\mathcal A$ be a finite set of
interaction factors.  Agent $i$ has a standard Borel state space
$(\mathsf Y_i,\mathscr Y_i)$.  Its state may contain a belief-law point in
$(\mathcal E_b)_{c_i}$, a model-law point in $(\mathcal E_m)_{c_i}$, and
declared auxiliary latent variables.  Write
$(\mathsf Y,\mathscr Y)=\prod_{i\in V}(\mathsf Y_i,\mathscr Y_i)$ and let
\begin{equation}
P_0:(\mathsf X,\mathscr X)\rightsquigarrow(\mathsf Y,\mathscr Y)
\label{eq:obs-context-baseline-kernel}
\end{equation}
be a normalized probability kernel.  Each factor $a$ has a finite scope
$\partial a\subseteq V$, a standard Borel record space
$(\mathsf O_a,\mathscr O_a)$, and a normalized Markov kernel
\begin{equation}
K_a:(\mathsf X\times\mathsf Y_{\partial a},
\mathscr X\otimes\mathscr Y_{\partial a})
\rightsquigarrow(\mathsf O_a,\mathscr O_a).
\label{eq:obs-interaction-kernel}
\end{equation}
For the density presentation used in this chapter, one fixed probability or
sigma-finite measure $\nu_a$ dominates the whole family
$\{K_a(X,y_{\partial a},\cdot):(X,y_{\partial a})
\in\mathsf X\times\mathsf Y_{\partial a}\}$, and a finite jointly measurable
version
\begin{equation}
\ell_a:\mathsf X\times\mathsf Y_{\partial a}\times\mathsf O_a
\longrightarrow[0,+\infty),
\qquad
K_a(X,y_{\partial a},do_a)
=\ell_a(o_a\given X,y_{\partial a})\nu_a(do_a)
\label{eq:obs-interaction-density}
\end{equation}
is fixed once and for all.  Thus
$\int\ell_a(o_a\given X,y_{\partial a})\nu_a(do_a)=1$ for every
$(X,y_{\partial a})$.  The common $\nu_a$ and the displayed joint
measurability are family-level hypotheses, not recordwise choices.
The kernels and $P_0$ are fixed by the generative parameters and structural
data; none reads a recognition law, a variational parameter, or a live
posterior.  Their local coordinate presentations are required to transform by
the common passive pushforward induced by a change of bundle frames, so the
underlying measures and record probabilities are frame independent.
\status{HYPOTHESIS}

This hypothesis permits correlated baseline states.  A graph-local Markov
blanket requires the additional condition that the regular conditionals of
$P_0$ have the corresponding locality.  Bundle covariance is imposed by
requiring each $K_a$ to be invariant under simultaneous changes of the
incident agent frames and the factor frame.  One concrete realization gives
factor $a$ its own principal frame and incidence-leg transports
$U^x_{a\leftarrow i}$ in each channel $x\in\{b,m\}$; the likelihood is then a
measurable function of the dressed variables
$\rho_x(U^x_{a\leftarrow i})y_i^x$.  The two representations remain distinct,
and a belief--model comparison still requires the declared cross morphisms of
\Cref{sec:geo-cross-morphisms}.  \status{DEFINITION}

Conditional independence of the records given the agent configuration defines
the joint law
\begin{equation}
P_\theta(dy,do\given X)
=P_0(dy\given X)\prod_{a\in\mathcal A}
K_a(X,y_{\partial a},do_a).
\label{eq:obs-interaction-joint}
\end{equation}

\propositionheading{Normalization on arbitrary interaction cycles}{prop:obs-interaction-normalization}
The law in \eqref{eq:obs-interaction-joint} is normalized for every finite
hypergraph, including cyclic ones.  \status{ESTABLISHED}

\paragraph{Proof.}
Integrate the record coordinates in any order.  Each normalized kernel
integrates to one, leaving the normalized baseline law $P_0(\cdot\given X)$.
The factor scopes
are permitted to overlap because the record variables are conditionally independent
children of one latent configuration.  This argument would not validate an
arbitrary cycle of separately prescribed agent conditionals, which need not be
compatible with any joint law.  $\square$

For the remainder of this chapter fix one $X\in\mathsf X$.  We suppress that
fixed argument consistently:
$P_0(dy):=P_0(dy\given X)$,
$K_a(y_{\partial a},do_a):=K_a(X,y_{\partial a},do_a)$, and
$\ell_a(o_a\given y_{\partial a}):=
\ell_a(o_a\given X,y_{\partial a})$.  Every descendant symbol
$L_o,H_o,Z,\Pi_o$ and its block analogue is relative to this same fixed
context unless an $X$ argument is restored explicitly.

With $\nu=\bigotimes_{a\in\mathcal A}\nu_a$, the complete record density is
the finite jointly measurable function
\begin{equation}
\ell(o\given y)=\prod_{a\in\mathcal A}
\ell_a(o_a\given y_{\partial a}),
\qquad
P_{\theta,X}(dy,do)=P_0(dy)\ell(o\given y)\nu(do).
\label{eq:obs-complete-record-kernel-density}
\end{equation}
For a fixed record $o$, define its likelihood before taking logarithms:
\begin{equation}
L_o(y):=\ell(o\given y)
=\prod_{a\in\mathcal A}\ell_a(o_a\given y_{\partial a})
\in[0,+\infty).
\label{eq:obs-complete-record-likelihood}
\end{equation}
Its extended interaction energy is
\begin{equation}
H_o(y):=-\log L_o(y)
=\sum_{a\in\mathcal A}-\log\ell_a(o_a\given y_{\partial a}),
\qquad -\log0:=+\infty.
\label{eq:obs-interaction-energy}
\end{equation}
Because each factor density is finite, $H_o$ never equals $-\infty$, although
individual negative log densities and their finite sum can be negative.  A
binary success record with
$K_a(1\given y)=\exp[-E_a(y)]$ is a valid special case only when
$E_a\in[0,+\infty]$.  General observations are represented by the normalized
kernels, not by treating an arbitrary exponential as a probability.
\status{ESTABLISHED}

\section{One collective VFE}
\label{sec:local-collective-vfe}

Fix a regular record $o$ for which
\begin{equation}
0<Z(o):=\int_{\mathsf Y}L_o(y)P_0(dy)<+\infty,
\label{eq:obs-global-evidence}
\end{equation}
and let
\begin{equation}
\Pi_o(dy)=Z(o)^{-1}L_o(y)P_0(dy)
\label{eq:obs-global-posterior}
\end{equation}
be the conditional posterior.  The regular-observation qualifications of
\Cref{ch:elbo} apply.

\theoremheading{Collective interaction VFE}{thm:obs-collective-vfe}
For every probability measure $Q$ on $\mathsf Y$, define the collective
extended ELBO and VFE by
\begin{equation}
\boxed{
\mathcal F_o^{\rm ext}(Q):=-\mathcal L_o^{\rm ext}(Q)
:=-\log Z(o)+\KL(Q\Vert\Pi_o).}
\label{eq:obs-collective-vfe}
\end{equation}
This takes values in $\mathbb R\cup\{+\infty\}$ and is the specialization of
\Cref{def:elbo-extended,thm:elbo-extended-gap}.  Equivalently,
$\mathcal L_o^{\rm ext}(Q)\leq\log Z(o)$ for every $Q$, with equality exactly
when $Q=\Pi_o$ as measures.  No domination or mean-field assumption on $Q$ is
needed.  \status{ESTABLISHED}

\paragraph{Proof.}
Apply the extended gap theorem \Cref{thm:elbo-extended-gap} with reference
$P_0$ and likelihood $L_o$.  Nonnegativity and the equality condition are the
corresponding properties of relative entropy.  $\square$

If $Q\ll P_0$ and the separate complexity and energy terms are defined without
an indeterminate $+\infty-\infty$, the same quantity has the familiar energy
presentation
\begin{equation}
\mathcal F_o^{\rm ext}(Q)
=\KL(Q\Vert P_0)+\E_QH_o.
\label{eq:obs-collective-energy-form}
\end{equation}
This is a derived representation, not the definition at singular $Q$.
\status{ESTABLISHED}

There is one term in $H_o$ for each actual interaction record.  An undirected
record shared by agents $i$ and $j$ occurs once.  Writing both $E_{ij}$ and
$E_{ji}$ in the global law asserts two directed records, not two descriptions
of one record.  \status{ESTABLISHED}

\section{The conditional VFE of an agent or meta-agent}
\label{sec:local-conditional-vfe}

Let $B\subseteq V$ be nonempty and write $b=y_{B^c}$.  The same construction
therefore treats a singleton agent $B=\{i\}$ and a declared meta-agent block
$B$ without changing probability models.  Define
\begin{equation}
\mathcal A_B=\{a\in\mathcal A:\partial a\cap B\ne\varnothing\},
\qquad
g_{B,o}(y_B;b)=\prod_{a\in\mathcal A_B}
\ell_a(o_a\given y_{\partial a}),
\label{eq:obs-block-incident-likelihood}
\end{equation}
and the likelihood of records whose scopes miss $B$,
\begin{equation}
L_{\bar B,o}(b)=
\prod_{a:\partial a\cap B=\varnothing}
\ell_a(o_a\given b_{\partial a}),
\qquad
L_o(y_B,b)=g_{B,o}(y_B;b)L_{\bar B,o}(b).
\label{eq:obs-block-outside-likelihood}
\end{equation}
The corresponding extended energies retain the useful additive notation
\begin{equation}
H_{B,o}(y_B;b)=-\log g_{B,o}(y_B;b)
=\sum_{a\in\mathcal A_B}-\log\ell_a(o_a\given y_{\partial a})
\label{eq:obs-block-incident-energy}
\end{equation}
and
\begin{equation}
H_{\bar B,o}(b)=-\log L_{\bar B,o}(b),
\qquad H_o(y)=H_{B,o}(y_B;b)+H_{\bar B,o}(b),
\label{eq:obs-block-outside-energy}
\end{equation}
as an identity in $\mathbb R\cup\{+\infty\}$.

Fix once and for all a measurable regular-conditional version
$b\mapsto P_{0,B}(\cdot\given b)$ in the disintegration
\begin{equation}
P_0(dy_B,db)=P_{0,B^c}(db)P_{0,B}(dy_B\given b).
\label{eq:obs-baseline-block-disintegration}
\end{equation}
Every pointwise formula in $b$ below is relative to this declared version.
Two baseline-conditional versions agree $P_{0,B^c}$-almost everywhere and
hence, by \eqref{eq:obs-posterior-outside}, $\Pi_{o,B^c}$-almost everywhere.
No pointwise version invariance is claimed beyond these almost-sure statements:
their values and all descendant formulas at exceptional $b$ are not determined
by the joint law.  Two different conditional likelihoods are needed.
The incident-factor evidence is
\begin{equation}
Z_B(b)=\int g_{B,o}(y_B;b)P_{0,B}(dy_B\given b),
\label{eq:obs-local-evidence}
\end{equation}
whereas the full likelihood of the outside state is defined directly by
\begin{equation}
w_B(b)=\int L_o(y_B,b)P_{0,B}(dy_B\given b).
\label{eq:obs-full-complement-likelihood}
\end{equation}
The direct definition is valid even where the formal product
$L_{\bar B,o}(b)Z_B(b)$ would have the indeterminate form $0\cdot\infty$.
Tonelli's theorem gives
\begin{equation}
Z(o)=\int w_B(b)P_{0,B^c}(db),
\qquad
\Pi_{o,B^c}(db)=Z(o)^{-1}w_B(b)P_{0,B^c}(db).
\label{eq:obs-posterior-outside}
\end{equation}
Let $\mathsf R_{B,o}$ be the measurable set on which
\begin{equation}
0<w_B(b)<+\infty,
\quad 0<L_{\bar B,o}(b)<+\infty,
\quad 0<Z_B(b)<+\infty,
\quad w_B(b)=L_{\bar B,o}(b)Z_B(b).
\label{eq:obs-block-regular-set}
\end{equation}
Then $\Pi_{o,B^c}(\mathsf R_{B,o})=1$.  Indeed, finite collective evidence
makes $w_B<+\infty$ outside a $P_{0,B^c}$-null set, and the posterior outside
is supported on $w_B>0$.  On that support the finite-factor hypothesis makes
$L_{\bar B,o}$ finite; positivity of $w_B$ forces
$L_{\bar B,o}>0$, after which ordinary factorization under the conditional
integral gives $Z_B=w_B/L_{\bar B,o}\in(0,+\infty)$.  The regular set need not
be $P_{0,B^c}$-full.  \status{ESTABLISHED}

For $b\in\mathsf R_{B,o}$ relative to the fixed baseline-conditional version,
the posterior conditional is
\begin{equation}
\Pi_{o,B}(dy_B\given b)
=Z_B(b)^{-1}g_{B,o}(y_B;b)P_{0,B}(dy_B\given b),
\label{eq:obs-local-posterior}
\end{equation}
and, after choosing any measurable probability-kernel extension off
$\mathsf R_{B,o}$ (for example the fixed baseline conditional there), it is a
version for which
\begin{equation}
\Pi_o(dy_B,db)=
\Pi_{o,B^c}(db)\Pi_{o,B}(dy_B\given b).
\label{eq:obs-posterior-block-disintegration}
\end{equation}
The posterior joint determines this kernel only
$\Pi_{o,B^c}$-almost everywhere.  In particular, neither the conditional nor
the associated pointwise local VFE at an outside-null $b$ is an invariant of
the joint law.  All almost-sure local/global statements below use
$Q_{B^c}\ll\Pi_{o,B^c}$ and are therefore independent of these choices.
\status{ESTABLISHED}

\theoremheading{Local multi-agent ELBO}{thm:obs-local-multiagent-elbo}
For $b\in\mathsf R_{B,o}$ and every probability measure $r_B(\cdot\given b)$,
define
\begin{equation}
\boxed{
\mathcal F_{B,o}^{\rm ext}(r_B;b)
:=-\mathcal L_{B,o}^{\rm ext}(r_B;b)
:=-\log Z_B(b)
+\KL\bigl(r_B(\cdot\given b)\Vert\Pi_{o,B}(\cdot\given b)\bigr).}
\label{eq:obs-local-vfe}
\end{equation}
Thus $\mathcal L_{B,o}^{\rm ext}(r_B;b)\leq\log Z_B(b)$, with equality
exactly at $r_B=\Pi_{o,B}$.  The singleton $B=\{i\}$ is one agent's local VFE;
an arbitrary nonempty $B$ is the same construction for a meta-agent.
No domination or mean-field condition on $r_B$ is needed.
\status{ESTABLISHED}

\paragraph{Proof.}
Apply \Cref{thm:elbo-extended-gap} conditionally with reference
$P_{0,B}(\cdot\given b)$ and likelihood $g_{B,o}(\cdot;b)$.  The equality
condition follows from vanishing relative entropy.  $\square$

On the dominated energy tier, when its two terms are separately defined
without an indeterminate infinity, the local functional is also
\begin{equation}
\mathcal F_{B,o}^{\rm ext}(r_B;b)
=\KL\bigl(r_B(\cdot\given b)\Vert P_{0,B}(\cdot\given b)\bigr)
+\E_{r_B(\cdot\given b)}H_{B,o}(Y_B;b).
\label{eq:obs-local-energy-form}
\end{equation}

The exact relative-entropy chain rule is most safely written against the
posterior.  If $Q(dy_B,db)=Q_{B^c}(db)r_B(dy_B\given b)$ and
$Q_{B^c}\ll\Pi_{o,B^c}$, then, in $[0,+\infty]$,
\begin{equation}
\KL(Q\Vert\Pi_o)
=\KL(Q_{B^c}\Vert\Pi_{o,B^c})
+\E_{Q_{B^c}}
\KL\bigl(r_B(\cdot\given Y_{B^c})
\Vert\Pi_{o,B}(\cdot\given Y_{B^c})\bigr).
\label{eq:obs-posterior-block-kl-chain}
\end{equation}
The baseline chain rule, whenever its regular conditional versions are used,
likewise reads
\begin{equation}
\KL(Q\Vert P_0)
=\KL(Q_{B^c}\Vert P_{0,B^c})
+\E_{Q_{B^c}}
\KL\bigl(r_B(\cdot\given Y_{B^c})
\Vert P_{0,B}(\cdot\given Y_{B^c})\bigr).
\label{eq:obs-block-kl-chain}
\end{equation}

\theoremheading{Exact local--global potential identity}{thm:obs-local-global-potential}
Let $Q=Q_{B^c}r_B$ and $Q'=Q_{B^c}r'_B$ have the same outside marginal with
$Q_{B^c}\ll\Pi_{o,B^c}$.  Suppose
$\KL(Q\Vert\Pi_o)<+\infty$ and
$\KL(Q'\Vert\Pi_o)<+\infty$; the chain rule then makes both conditional-KL
fields integrable, so the following subtraction is finite.  Then
\begin{equation}
\mathcal F_o^{\rm ext}(Q')-\mathcal F_o^{\rm ext}(Q)
=\E_{Q_{B^c}}\left[
\mathcal F_{B,o}^{\rm ext}(r'_B;Y_{B^c})
-\mathcal F_{B,o}^{\rm ext}(r_B;Y_{B^c})
\right].
\label{eq:obs-local-global-potential}
\end{equation}
Thus every exact local conditional update is a coordinate update of the same
collective VFE.  \status{ESTABLISHED}

\paragraph{Proof.}
Apply \eqref{eq:obs-posterior-block-kl-chain} to both laws.  Their outside KL
terms agree.  The common finite constant $-\log Z(o)$ cancels, as does the
pointwise term $-\log Z_B(Y_{B^c})$ inside the two local functionals.  The
finite-difference hypotheses permit the remaining subtraction under the
integral.  $\square$

When all baseline KL terms and positive and negative energy parts are
integrable, the same result admits the familiar finite decomposition
\begin{equation}
\begin{split}
\mathcal F_o^{\rm ext}(Q)
&=\KL(Q_{B^c}\Vert P_{0,B^c})\\
&\quad+\E_{Q_{B^c}}\left[
H_{\bar B,o}(Y_{B^c})
+\mathcal F_{B,o}^{\rm ext}(r_B;Y_{B^c})
\right].
\end{split}
\label{eq:obs-local-global-decomposition}
\end{equation}
This theorem licenses sequential exact coordinate minimization, damped updates
accepted by the collective objective, and continuous block-gradient dynamics.
Independently replacing all correlated full conditionals need not define any
joint recognition law, so such a parallel prescription is not licensed.
\status{ESTABLISHED}

\begin{figure}[htbp]
\centering
\begin{tikzpicture}[>=stealth]
  \node[draw,circle,fill=lightblue,minimum size=0.85cm] (a1) at (0,1.5) {$1$};
  \node[draw,circle,fill=lightblue,minimum size=0.85cm] (a2) at (2.1,2.3) {$2$};
  \node[draw,circle,fill=lightgreen,minimum size=0.85cm] (a3) at (4.2,1.5) {$3$};
  \node[draw,circle,fill=lightgreen,minimum size=0.85cm] (a4) at (5.2,-0.2) {$4$};
  \node[draw,circle,fill=lightblue,minimum size=0.85cm] (a5) at (1.2,-0.5) {$5$};
  \node[draw,rectangle,fill=lightorange,minimum size=0.52cm] (f12) at (1.05,1.9) {$o_a$};
  \node[draw,rectangle,fill=lightorange,minimum size=0.52cm] (f235) at (2.65,0.45) {$o_b$};
  \node[draw,rectangle,fill=lightorange,minimum size=0.52cm] (f34) at (4.7,0.65) {$o_c$};
  \draw (a1)--(f12)--(a2);
  \draw (a2)--(f235)--(a3);
  \draw (a5)--(f235);
  \draw (a3)--(f34)--(a4);
  \draw[rounded corners,very thick,sciencegreen]
    (3.55,2.15) rectangle (5.75,-0.85);
  \node[text=sciencegreen] at (4.65,2.48) {block $B=\{3,4\}$};
  \draw[->,very thick,primaryblue] (6.05,0.65)--(7.35,0.65);
  \node[align=center,draw,rounded corners,fill=lightblue,text width=3.5cm]
    at (9.25,0.65)
    {$\Delta\mathcal F_o^{\rm ext}
      =\E_{Q_{B^c}}\Delta\mathcal F_{B,o}^{\rm ext}$\\
     incident records: $o_b,o_c$};
\end{tikzpicture}
\caption{A block-local VFE contains every factor incident to the block.  The
factor $o_b$ crosses the boundary and therefore enters the block conditional;
$o_a$ is outside and cancels.  The equality is an exact-potential identity,
not a sum of overlapping local objectives.}
\label{fig:obs-local-global-potential}
\end{figure}

\section{Additive accounting is a different construction}
\label{sec:local-additive-accounting}

When $P_0=\bigotimes_i\rho_i$, define the total correlation of an arbitrary
joint recognition law $Q$ by
\begin{equation}
\operatorname{TC}(Q)=\KL\left(Q\middle\Vert\bigotimes_iQ_i\right).
\label{eq:obs-total-correlation}
\end{equation}
Without assuming mean-field factorization, the extended chain rule gives
\begin{equation}
\KL(Q\Vert P_0)
=\operatorname{TC}(Q)+\sum_i\KL(Q_i\Vert\rho_i)
\quad\text{in }[0,+\infty].
\label{eq:obs-correlated-complexity-ledger}
\end{equation}
Consequently, whenever the energy presentation
\eqref{eq:obs-collective-energy-form} is defined without an indeterminate
infinity,
\begin{equation}
\mathcal F_o^{\rm ext}(Q)
=\operatorname{TC}(Q)
+\sum_i\KL(Q_i\Vert\rho_i)
+\E_Q\left[\sum_{a\in\mathcal A}E_{a,o}\right],
\qquad E_{a,o}=-\log\ell_a(o_a\given Y_{\partial a}).
\label{eq:obs-global-ledger}
\end{equation}
\status{ESTABLISHED}

Equation~\eqref{eq:obs-global-ledger} is an additive ledger: each interaction
appears once inside the pointwise extended factor sum.  It can be rewritten as
$\sum_a\E_QE_{a,o}$ when every $E_{a,o}\in L^1(Q)$, or under another
condition excluding opposing infinities.  Moreover,
$\operatorname{TC}(Q)=0$ exactly when $Q=\bigotimes_iQ_i$; mean field is thus a
special case of the ledger, not a premise of it.

For the singleton conditional objectives actually defined above, assume the
declared interaction factors have nonempty scopes.  Their incident energies
then obey the exact pointwise identity
\begin{equation}
H_{\{i\},o}(y_i;y_{-i})
=\sum_{a:i\in\partial a}E_{a,o}(y_{\partial a}),
\qquad
\sum_{i\in V}H_{\{i\},o}(y_i;y_{-i})
=\sum_{a\in\mathcal A}|\partial a|E_{a,o}(y_{\partial a}).
\label{eq:obs-singleton-incident-counting}
\end{equation}
An empty-scope record, if separately permitted, is instead a label-independent
global constant, and the right-hand sum must then be restricted to
$|\partial a|>0$.  Because every
$E_{a,o}$ lies in
$\mathbb R\cup\{+\infty\}$ and both index sets are finite, the rearrangement
is valid as an extended pointwise identity.  Thus the specific local VFEs of
\Cref{thm:obs-local-multiagent-elbo} have overlapping incident-energy
components and are not themselves the single-count ledger
\eqref{eq:obs-global-ledger}.  This is a counting identity, not a universal
no-go for arbitrary decompositions.

Ordered chain rules, a separate residual account, or nonlocal counting-number
shares---including compensating negative shares assigned away from an incident
coordinate---can restore an additive sum.  Such a construction depends on
extra bookkeeping and is not the conditional local VFE of
\Cref{thm:obs-local-multiagent-elbo}.
\status{ESTABLISHED}

\section{Attention as a latent interaction label}
\label{sec:local-attention-elbo}

Attention enters the fixed joint by a latent source label, not by allowing the
generative law to inspect recognition beliefs.  Let
$\mathcal I_{\rm att}\subseteq V$ be the finite set of receiving agents.  For
each $i\in\mathcal I_{\rm att}$, let $J_i$ augment the baseline independently
as
\begin{equation}
P_0^{\rm aug}(dy,dj)=P_0^Y(dy)
\prod_{i\in\mathcal I_{\rm att}}\pi_i(dj_i).
\label{eq:obs-attention-augmented-baseline}
\end{equation}
Suppose $\pi_{ij}>0$ on the finite declared source set and that every
$D_{ij}:\mathsf Y\to\mathbb R$ is finite and measurable.  At the selected
regular interaction record, suppose its density has source-dependent energy
\begin{equation}
\ell_i(o_i\given y,J_i=j)
=c_i(o_i,y)\exp[-D_{ij}(y)/\tau_i],
\qquad \tau_i>0,
\label{eq:obs-attention-likelihood}
\end{equation}
where $c_i$ is independent of $j$ and the normalized observation kernel is
part of the model data.  The energy $D_{ij}$ is a fixed gauge-invariant
function of latent sample-level states and declared transports.  It does not
contain a live recognition distribution.  Equation~\eqref{eq:obs-attention-likelihood}
is an identity on the selected observation slice; normalization of the kernel
and its values away from that slice remain part of the declared generative
model.  More strongly, at the complete selected record the full augmented
likelihood is required to factor as
\begin{equation}
L_o^{\rm aug}(y,j)
=L_o^Y(y)\prod_{i\in\mathcal I_{\rm att}}
c_i(o_i,y)\exp[-D_{ij_i}(y)/\tau_i],
\label{eq:obs-attention-augmented-likelihood}
\end{equation}
where $L_o^Y:\mathsf Y\to[0,+\infty)$ is a finite measurable function
independent of every source label and no other record or factor in the
generative law reads any $J_i$.  Equation
\eqref{eq:obs-attention-augmented-likelihood} is a label-exclusivity
hypothesis on the selected likelihood slice; normalization of the augmented
observation kernel over all possible records remains declared model data.
Without exclusivity, an additional $J_i$-dependent likelihood factor would
enter Bayes' formula and generally change the displayed softmax below.
\status{HYPOTHESIS}

The constant-row recognition calculation below additionally restricts the
augmented recognition law to the full factorization
\begin{equation}
Q(dy,dj)=Q_Y(dy)
\prod_{i\in\mathcal I_{\rm att}}\beta_i^Q(dj_i),
\qquad \beta_i^Q(dj_i)\ \text{independent of }y.
\label{eq:obs-attention-recognition-factorization}
\end{equation}
This assumption imposes both conditional independence among the source labels
and independence of every label row from the latent state.  It is not implied
by the generative law.  For the finite-valued formulas below, also assume
$\E_{Q_Y}|D_{ij}|<+\infty$ on every active source and
$\E_{Q_Y}|\log c_i(o_i,Y)|<+\infty$.  \status{HYPOTHESIS}

\propositionheading{Exact posterior and recognition attention}{prop:obs-attention-elbo}
Under the generative label-exclusivity hypothesis
\eqref{eq:obs-attention-augmented-likelihood}, conditioning on $y$ and the
complete interaction record gives, on the positive-likelihood set
$L_o^Y(y)\prod_{k\in\mathcal I_{\rm att}}c_k(o_k,y)>0$, the generative
posterior row
\begin{equation}
\beta^P_{ij}(y)
=\frac{\pi_{ij}\exp[-D_{ij}(y)/\tau_i]}
{\sum_k\pi_{ik}\exp[-D_{ik}(y)/\tau_i]}.
\label{eq:obs-attention-posterior}
\end{equation}
and the same display may be chosen as its version off that set.  Under the
additional recognition restriction
\eqref{eq:obs-attention-recognition-factorization}, the exact categorical
contribution of row $i$ to the collective VFE is
\begin{equation}
-\E_{Q_Y}\log c_i(o_i,Y)+\Fenergy_i^{\rm att}(\beta_i^Q),
\label{eq:obs-attention-full-contribution}
\end{equation}
where the first term is independent of $\beta_i^Q$ and
\begin{equation}
\Fenergy_i^{\rm att}(\beta^Q_i)
=\KL(\beta^Q_i\Vert\pi_i)
+\frac{1}{\tau_i}\sum_j\beta^Q_{ij}\E_{Q_Y}D_{ij}.
\label{eq:obs-attention-vfe}
\end{equation}
Its unique interior minimizer on the positive-prior source simplex is
\begin{equation}
\beta^{Q\star}_{ij}
=\frac{\pi_{ij}\exp[-\E_{Q_Y}D_{ij}/\tau_i]}
{\sum_k\pi_{ik}\exp[-\E_{Q_Y}D_{ik}/\tau_i]}.
\label{eq:obs-attention-recognition-optimum}
\end{equation}
\status{ESTABLISHED}

\paragraph{Proof.}
The product label prior and the exclusive likelihood
\eqref{eq:obs-attention-augmented-likelihood} make the conditional posterior
over $J$ factor by rows; every factor not containing $j_i$ cancels.  Bayes'
formula then gives \eqref{eq:obs-attention-posterior}.  With
$Z_i^{\rm att}(y)=\sum_k\pi_{ik}\exp[-D_{ik}(y)/\tau_i]$, expanding the
categorical relative entropy to that posterior gives
\[
\KL(\beta^Q_i\Vert\beta^P_i)
=\KL(\beta^Q_i\Vert\pi_i)
+\tau_i^{-1}\sum_j\beta^Q_{ij}D_{ij}+\log Z_i^{\rm att}(y).
\]
Subtracting $\log Z_i^{\rm att}(y)$ and the source-independent log likelihood
$\log c_i(o_i,y)$, then taking the recognition expectation, proves
\eqref{eq:obs-attention-full-contribution} and
\eqref{eq:obs-attention-vfe}.
Lagrange multiplication under $\sum_j\beta^Q_{ij}=1$ proves
\eqref{eq:obs-attention-recognition-optimum}.  $\square$

The two softmax rows answer different variational questions.
$\beta^P_i(y)$ is the generative posterior at a fixed latent state, whereas
$\beta_i^{Q\star}$ is the optimum after forcing one recognition row to be
constant in $y$.  In general
$\beta_i^{Q\star}\ne\E_{Q_Y}[\beta_i^P(Y)]$; the two formulas agree pointwise
when all differences $D_{ij}(Y)-D_{ik}(Y)$ are $Q_Y$-almost surely constant.
\status{ESTABLISHED}

For completeness, still under
\eqref{eq:obs-attention-augmented-likelihood}, let an arbitrary recognition law
disintegrate as
$Q(dy,dj)=Q_Y(dy)Q_{J\mid Y}(dj\given y)$.  With
$Q_{J_i\mid Y}$ denoting its conditional row marginals and the label-free term
$-\E_{Q_Y}\log L_o^Y(Y)$ accounted for in the state sector, the exact
attention-label contribution is
\begin{equation}
\begin{split}
&-\sum_{i\in\mathcal I_{\rm att}}\E_{Q_Y}\log c_i(o_i,Y)
+\E_{Q_Y}\operatorname{TC}(Q_{J\mid Y})\\
&\quad+\sum_{i\in\mathcal I_{\rm att}}\E_{Q_Y}\left[
\KL(Q_{J_i\mid Y}\Vert\pi_i)
+\frac{1}{\tau_i}
\E_{Q_{J_i\mid Y}}D_{iJ_i}(Y)
\right],
\end{split}
\label{eq:obs-attention-correlated-ledger}
\end{equation}
where
$\operatorname{TC}(Q_{J\mid Y=y})=
\KL\left(Q_{J\mid Y=y}\middle\Vert
\prod_{i\in\mathcal I_{\rm att}}Q_{J_i\mid Y=y}\right)$.
Thus merely requiring every row marginal to be independent of $y$ does not
justify \eqref{eq:obs-attention-vfe}: a correlated conditional law retains the
nonnegative conditional-total-correlation term.  If conditional dependence is
unrestricted, minimizing over all $Q_{J\mid Y=y}$ gives the product of the
pointwise rows \eqref{eq:obs-attention-posterior}.  Under a constrained coupled
family, the coupling term or constraint enters the coordinate optimum, so a
row softmax is not asserted.  \status{ESTABLISHED}

The familiar row functional
$\sum_j\beta_{ij}\E_{Q_Y}D_{ij}+\tau_i\KL(\beta_i\Vert\pi_i)$ is
$\tau_i\Fenergy_i^{\rm att}$.  It has the same row minimizer, but for
$\tau_i\ne1$ it is not an independently weighted sector of one standard
global ELBO unless the entire objective is scaled or a different generalized
free energy is declared.  \status{ESTABLISHED}

Holding the latent recognition law $Q_Y$, and hence every
$\E_{Q_Y}D_{ij}$, fixed
during this row-coordinate flow, the Fisher natural gradient of
\eqref{eq:obs-attention-vfe} on the open simplex is the replicator equation
\begin{equation}
\dot\beta^Q_{ij}
=-\gamma_i\beta^Q_{ij}
\left(c_{ij}-\sum_k\beta^Q_{ik}c_{ik}\right),
\qquad
c_{ij}=\log\frac{\beta^Q_{ij}}{\pi_{ij}}
+\frac{\E_{Q_Y}D_{ij}}{\tau_i}.
\label{eq:obs-attention-replicator}
\end{equation}
for $\gamma_i\geq0$.  Its partial contribution to the row-free-energy rate is
\begin{equation}
\left.\frac{d}{dt}\Fenergy_i^{\rm att}\right|_{\beta_i}
=-\gamma_i\operatorname{Var}_{\beta^Q_i}(c_{ij})\leq0.
\label{eq:obs-attention-dissipation}
\end{equation}
If $Q_Y$ evolves simultaneously, its additional chain-rule term must be
included.  \status{ESTABLISHED}

\section{Local evolution is collective descent}
\label{sec:local-natural-gradient}

Let a regular product recognition family $Q_\eta$ on an open parameter domain
have coordinates $\eta=(\eta_i)_i$ and block-diagonal Fisher metric
$G(\eta)=\bigoplus_iG_i(\eta_i)$, including categorical blocks when present.
Assume $F(\eta):=\mathcal F_o^{\rm ext}(Q_\eta)\in\mathbb R$ is $C^1$, every
$G_i(\eta_i)$ is positive definite, every $\gamma_i\geq0$, and the trajectory
remains in this domain.
By \Cref{thm:obs-local-global-potential}, the gradient of the outside-averaged
local VFE is the corresponding block gradient of $F$.  Hence
\begin{equation}
\dot\eta_i=-\gamma_iG_i^{-1}\nabla_{\eta_i}F
\label{eq:obs-local-natural-gradient}
\end{equation}
obeys
\begin{equation}
\frac{d}{dt}F
=-\sum_i\gamma_i
(\nabla_{\eta_i}F)^\top
G_i^{-1}(\nabla_{\eta_i}F)\leq0.
\label{eq:obs-global-dissipation}
\end{equation}
Thus an agent following its exact local VFE follows a block of the collective
VFE flow.  \status{ESTABLISHED}

Block orthogonality is load bearing.  With a nondiagonal Fisher metric, the
global natural gradient mixes agents and independent local inversions are a
different dynamics.  Hard support boundaries and finite steps also require a
separate projected flow, line search, damping, or acceptance test.
\status{ESTABLISHED}

The variable \(t\) in \eqref{eq:obs-local-natural-gradient} parameterizes the
declared block flow.  Passing from a regular solution curve to its oriented
unparameterized orbit, and then measuring that orbit by Fisher length, are the
separate constructions of
\Cref{def:hist-oriented-history,thm:hist-fisher-clock-invariance}.  A common
positive scalar mobility changes only the parameterization, whereas unequal
block rates \(\gamma_i\) generally change the orbit.  \status{ESTABLISHED}

No physical-time interpretation of \(t\) or of the resulting Fisher duration
is asserted.  \status{NOT-CLAIMED}

\section{Are observations only agent--agent interactions?}
\label{sec:local-observation-equivalence}

\definitionheading{Operational environment node}{def:obs-operational-environment-node}
An operational environment node is a standard Borel state space with a
normalized state law and a measurable message kernel.  This definition does
not assert autonomy, self-modeling, or biological agency.  \status{DEFINITION}

\theoremheading{Observation--interaction kernel equivalence}{thm:obs-agent-interaction-equivalence}
Every normalized standard-Borel observation kernel $K(do\given y)$ can be
realized as a message from an operational environment node: there are
$U\sim\operatorname{Unif}[0,1]$ and a measurable map $F$ such that
\begin{equation}
O=F(Y,U),
\qquad
\Pr(F(y,U)\in A)=K(A\given y).
\label{eq:obs-randomization}
\end{equation}
Conversely, marginalizing any environment state and message policy produces a
normalized observation kernel.  \status{ESTABLISHED}

\paragraph{Proof.}
The forward statement is the randomization lemma for kernels on standard
Borel spaces \citep{Kallenberg2021}.  For the converse, if the environment has
law $\rho(du)$ and policy $M(y,u,A)$, then
$K(A\given y)=\int M(y,u,A)\rho(du)$ is measurable in $y$, countably additive
in $A$, and normalized.  $\square$

The theorem gives an exact operationally agent-only presentation when every
datum is represented as a clamped state or message of a boundary node and all
exogenous noise is included among those nodes.  It does not eliminate the
conditioning role.  If $O$ denotes the complete record map, the ELBO is
conditioned on the retained sigma-algebra $\mathscr O=\sigma(O)$.  Replacing
the environment by boundary nodes preserves the same random variable $O$ and
hence the same $\mathscr O$; it does not delete it.  In particular,
\begin{equation}
I(Y;O)=\int
\KL\bigl(P(dY\given o)\Vert P^Y(dY)\bigr)P^O(do),
\label{eq:obs-conditioning-information}
\end{equation}
with values in $[0,+\infty]$.  This quantity measures how much the retained
interaction record changes the latent law on average.  \status{ESTABLISHED}

More generally, retaining only a statistic $S=T(O)$ gives the posterior
$P(dY\given S)$, not $P(dY\given O)$.  The deterministic-statistic chain rule
\begin{equation}
I(Y;O)=I(Y;S)+I(Y;O\given S)
\label{eq:obs-conditioning-statistic-chain}
\end{equation}
shows that this replacement is posterior-sufficient exactly when
$I(Y;O\given S)=0$ (up to the usual almost-sure choice of conditional
versions).  Complete deletion makes $S$ constant and leaves $P^Y$.  A finite
witness is $Y\sim\operatorname{Bernoulli}(1/2)$ with the deterministic message
$O=Y$: then $I(Y;O)=\log2$, $P(dY\given O=o)=\delta_o$, while deletion returns
the nondegenerate prior.  Thus an agent-only redescription cannot remove the
conditioning sigma-algebra unless it also accepts the corresponding loss of
information.  \status{ESTABLISHED}

Nor does the operational equivalence prove that a random seed is an autonomous
agent.  If physical agency requires persistent state, a Markov blanket,
action, or its own local VFE, those are additional hypotheses.  An
$Y$-independent kernel is a boundary message with zero mutual information and
therefore a null interaction in the informational sense.
\status{ESTABLISHED}

The resulting answer is two-tiered.  The probability theory can be written
entirely as interactions among agent and environment nodes, while the ELBO
still conditions on the sigma-algebra generated by those interaction records.
Whether every environment node deserves the stronger ontological name
``agent'' is not determined by the principal bundle or by variational
inference.  \status{NOT-CLAIMED}
